\documentclass[12pt]{beamer}
\usepackage{ceai-beamer}
\usepackage{ragged2e}
\AtBeginEnvironment{frame}{\justifying}
\AtBeginSection[]{}
\AtBeginSubsection[]{}

\title[Topic 18: AI Futures]{CE 488/5588: Applied AI\\Topic 18: AI Futures, World Models, Embodied Intelligence, and Engineering Judgment}
\author{}
\date{}

\newcommand{\source}[1]{\vspace{0.4em}{\tiny\color{gray}Source: #1\par}}

\begin{document}

\begin{frame}
  \titlepage
\end{frame}

\begin{frame}{Lecture Roadmap}
  \tableofcontents
\end{frame}

\section{From Topic 17 to Topic 18}

\begin{frame}{From Agentic Workflows to Frontier Questions}
  Previously, we examined why foundation models became useful (Topic 15), how retrieval and tools extend them (Topic 16), and what changes when those components become persistent workflows (Topic 17).

  In this chapter, we ask a different question:
  \begin{itemize}
    \item if current LLM systems are powerful but incomplete,
    \item what kinds of ideas are most likely to matter next,
    \item and what should engineers believe without drifting into mythology?
  \end{itemize}
\end{frame}

\begin{frame}{Why This Chapter Matters}
  The future-AI conversation is often distorted by false choices:
  \begin{itemize}
    \item either LLMs already contain everything needed,
    \item or one replacement paradigm will soon sweep them away,
    \item or engineering judgment will dissolve into prompting.
  \end{itemize}

  \begin{warningbox}{This chapter takes a stricter view}
    The most credible future is not one universal architecture. It is a layered stack in which different mechanisms handle interface, explicit structure, prediction, physical action, and oversight.
  \end{warningbox}
\end{frame}

\begin{frame}{Learning Objectives}
  After this lecture, students should be able to:
  \begin{itemize}
    \item explain why today's AI is best understood as a system stack rather than a standalone model,
    \item distinguish among foundation models, neuro-symbolic methods, world models, and embodied intelligence,
    \item describe why world models matter most when prediction and control are central,
    \item explain why embodiment matters for physical agency but not equally for every useful AI system,
    \item evaluate AGI and ASI claims using engineering standards of evidence, and
    \item explain why engineers remain responsible for constraints, verification, and accountability.
  \end{itemize}
\end{frame}

\section{The Current Paradigm}

\begin{frame}{The Current Paradigm: Foundation-Model Pragmatism}
  The foundation-model era has already changed coding, search, content generation, and parts of engineering workflow.

  In practice, useful systems are no longer just ``a model.'' They are models connected to:
  \begin{itemize}
    \item retrieval,
    \item memory,
    \item tools,
    \item environments, and
    \item human supervision.
  \end{itemize}

  \source{\href{https://arxiv.org/abs/2108.07258}{Bommasani et al. (2021): Foundation Models}}
\end{frame}

\begin{frame}{Why System Design Matters So Much}
  Current LLM systems are strong at:
  \begin{itemize}
    \item language interface,
    \item pattern extraction,
    \item summarization,
    \item procedural decomposition.
  \end{itemize}

  They are weaker at:
  \begin{itemize}
    \item robust grounding,
    \item stable physical state tracking,
    \item reasoning about open-ended physical consequence.
  \end{itemize}

  \source{\href{https://arxiv.org/abs/2310.07018}{Wang et al. (2023): NEWTON}; \href{https://arxiv.org/abs/2504.16074}{Qiu et al. (2025): PHYBench}}
\end{frame}

\begin{frame}{The Future AI Stack as an Engineering Stack}
  \begin{figure}
    \centering
    \begin{tikzpicture}[>=Stealth, font=\small, node distance=0.45cm, scale=0.90, transform shape]
      \tikzset{layer/.style={draw, rounded corners, minimum width=10.7cm, minimum height=0.9cm, align=center}}
      \node[layer, fill=blue!8] (fm) {Foundation model interface layer\\language, code, multimodal priors};
      \node[layer, fill=green!10, below=of fm] (tool) {System scaffolding\\retrieval, memory, tools, agent loops, verification};
      \node[layer, fill=orange!12, below=of tool] (world) {Task-specific intelligence layers\\constraints, world models, planners, domain simulators};
      \node[layer, fill=red!9, below=of world] (deploy) {Deployment reality\\physical environments, institutions, operators, accountability};
      \draw[->, thick] (fm) -- (tool);
      \draw[->, thick] (tool) -- (world);
      \draw[->, thick] (world) -- (deploy);
    \end{tikzpicture}
    \caption{Recreated teaching figure: a layered view of likely future AI system design.}
  \end{figure}
\end{frame}

\begin{frame}{Interpreting the Stack}
  \begin{conceptbox}{Main claim}
    The credible future is not ``LLMs disappear'' or ``LLMs become everything.'' Foundation models remain useful interface layers while other mechanisms become more important around them.
  \end{conceptbox}

  From an engineering perspective, the system boundary now includes:
  \begin{itemize}
    \item prompt structure,
    \item memory and retrieval,
    \item tool interfaces,
    \item verification layers,
    \item operators and institutions.
  \end{itemize}
\end{frame}

\begin{frame}{Why Language-First Systems Are Still Incomplete}
  A fluent model can imitate reasoning traces and summarize rules, but that is not the same as:
  \begin{itemize}
    \item keeping a stable world state,
    \item predicting physical consequences under intervention,
    \item recovering from novel physical failure,
    \item acting safely in real environments.
  \end{itemize}

  \begin{warningbox}{Engineering consequence}
    If the task involves explicit rules, physical prediction, or high consequence, language fluency alone is rarely enough.
  \end{warningbox}
\end{frame}

\section{Neural-Symbolic AI}

\begin{frame}{What Neural-Symbolic AI Means}
  \begin{conceptbox}{Definition}
    Neuro-symbolic AI combines learned subsymbolic representations with explicit symbolic machinery: rules, variables, programs, constraints, logic, or probabilistic reasoning.
  \end{conceptbox}

  Motivation:
  \begin{itemize}
    \item neural systems scale well in noisy, high-dimensional perception,
    \item symbolic systems preserve explicit structure and traceable inference,
    \item each side is strong where the other is weak.
  \end{itemize}

  \source{\href{https://arxiv.org/abs/1711.03902}{Besold et al. (2017)}; \href{https://arxiv.org/abs/2501.05435}{Colelough \& Regli (2025)}}
\end{frame}

\begin{frame}{Why Explicit Structure Keeps Returning}
  Engineers care about domains where the structure is not optional:
  \begin{itemize}
    \item laws,
    \item design constraints,
    \item operating procedures,
    \item ontologies,
    \item program semantics,
    \item scientific priors.
  \end{itemize}

  In such settings, symbolic structure improves:
  \begin{itemize}
    \item auditability,
    \item checkability,
    \item data efficiency,
    \item controllability.
  \end{itemize}
\end{frame}

\begin{frame}{Representative Neuro-Symbolic Systems}
  Representative systems make the idea concrete:
  \begin{itemize}
    \item \textbf{DeepProbLog}: learned perception inside probabilistic logic,
    \item \textbf{Neural Logic Machines}: relational and logical inductive bias,
    \item \textbf{NS-CL}: object-centric perception + semantic parsing + executable programs.
  \end{itemize}

  \source{\href{https://doi.org/10.1016/j.artint.2021.103504}{Manhaeve et al. (DeepProbLog)}; \href{https://arxiv.org/abs/1904.11694}{Dong et al. (NLM)}; \href{https://arxiv.org/abs/1904.12584}{Mao et al. (NS-CL)}}
\end{frame}

\begin{frame}{The Modern Extension: Agentic AI as System-Level Hybridization}
  Topic 17 described modern agentic systems in terms of:
  \begin{itemize}
    \item skills,
    \item tool interfaces,
    \item orchestration logic,
    \item explicit constraints,
    \item human review checkpoints.
  \end{itemize}

  \begin{conceptbox}{Careful claim}
    Many agentic systems are neuro-symbolic at the \emph{system level}, even when their neural core is still a large language model rather than a classical theorem prover.
  \end{conceptbox}
\end{frame}

\begin{frame}{Why Skills and Constraints Function Symbolically}
  When an engineer defines a reusable workflow, the engineer is not merely ``prompting'' a model.

  The engineer is imposing symbolic structure on the task:
  \begin{itemize}
    \item what steps exist,
    \item which tools may be called,
    \item what counts as acceptable evidence,
    \item when the system must stop,
    \item when a human must intervene.
  \end{itemize}

  If those rules can be inspected and revised by humans, they are functioning like a symbolic layer around a neural substrate.
\end{frame}

\begin{frame}{Why Natural Language Can Still Be Symbolic}
  Natural language carries symbolic structure when it is used to define:
  \begin{itemize}
    \item roles,
    \item branches,
    \item pass/fail criteria,
    \item escalation rules,
    \item scenario-specific constraints.
  \end{itemize}

  \begin{warningbox}{Important caution}
    The careful claim is not that every agentic system is a classical symbolic reasoner. The claim is that many useful agentic systems are neuro-symbolic in practice because neural models operate inside human-defined symbolic workflows.
  \end{warningbox}
\end{frame}

\begin{frame}{Engineering Example: Design Code Checking}
  \begin{examplebox}{Where this matters in engineering}
    Imagine an AI assistant for design code checking. A purely neural system may summarize text well, but the workflow still needs explicit clauses, pass/fail constraints, traceable justification, and human review of edge cases.
  \end{examplebox}

  That is exactly the kind of setting in which symbolic structure is not nostalgia. It is operationally useful.
\end{frame}

\section{World Models}

\begin{frame}{What Counts as a World Model?}
  The phrase \textbf{world model} is used in more than one way.

  \begin{itemize}
    \item In model-based reinforcement learning: a learned dynamics model that lets an agent imagine futures and plan in latent space.
    \item In the JEPA line: predictive latent representations that capture what is predictable and task-relevant without reconstructing every surface detail.
  \end{itemize}

  These communities overlap, but they should not be collapsed into a single slogan.
\end{frame}

\begin{frame}{The PlaNet and Dreamer Line}
  The model-based RL trajectory is already substantial:
  \begin{itemize}
    \item \textbf{PlaNet}: latent dynamics planning from pixels,
    \item \textbf{Dreamer}: behavior learned through latent imagination,
    \item \textbf{DreamerV3}: a much more scalable control program across many tasks.
  \end{itemize}

  This is not hype. It is a real research program with repeated empirical gains.

  \source{\href{https://arxiv.org/abs/1811.04551}{Hafner et al. (PlaNet)}}
\end{frame}

\begin{frame}{What Dreamer Adds Conceptually}
  Dreamer matters because it sharpens the idea of a world model:
  \begin{itemize}
    \item learn a latent state,
    \item imagine future trajectories in that latent space,
    \item optimize behavior using imagined rollouts rather than only expensive environment interaction.
  \end{itemize}

  \source{\href{https://arxiv.org/abs/1912.01603}{Hafner et al. (Dreamer)}}
\end{frame}

\begin{frame}{LeCun's Critique of Pure Next-Token Prediction}
  LeCun's roadmap argues that next-token prediction alone may be an insufficient route to autonomous intelligence because intelligence also requires:
  \begin{itemize}
    \item hierarchical representation,
    \item prediction,
    \item planning,
    \item action over latent world state.
  \end{itemize}

  \source{\href{https://openreview.net/forum?id=BZ5a1r-kVsf}{LeCun (2022): A Path Towards Autonomous Machine Intelligence}}
\end{frame}

\begin{frame}{I-JEPA and V-JEPA}
  I-JEPA and V-JEPA should be read as attempts to operationalize that roadmap:
  \begin{itemize}
    \item predict latent representations rather than reconstruct every pixel,
    \item emphasize what is predictable and semantically useful,
    \item move from image prediction to video and, increasingly, planning-relevant settings.
  \end{itemize}

  \source{\href{https://arxiv.org/abs/2301.08243}{Assran et al. (I-JEPA)}; \href{https://arxiv.org/abs/2404.08471}{Bardes et al. (V-JEPA)}}
\end{frame}

\begin{frame}{Autoregression versus World Models}
  \begin{figure}
    \centering
    \begin{tikzpicture}[>=Stealth, font=\small, node distance=0.6cm, scale=0.92, transform shape]
      \tikzset{box/.style={draw, rounded corners, minimum width=3.1cm, minimum height=0.9cm, align=center}}
      \node[box, fill=blue!8] (obs1) at (0,0) {Observation / tokens};
      \node[box, fill=blue!8] (pred1) at (3.9,0) {Next-step prediction};
      \node[box, fill=blue!8] (out1) at (7.8,0) {Generated continuation};
      \draw[->, thick] (obs1) -- (pred1);
      \draw[->, thick] (pred1) -- (out1);
      \node[box, fill=orange!10] (obs2) at (0,-2.2) {Observations + actions};
      \node[box, fill=orange!10] (state) at (3.9,-2.2) {Latent state / world model};
      \node[box, fill=orange!10] (plan) at (7.8,-2.2) {Imagined futures + planning};
      \draw[->, thick] (obs2) -- (state);
      \draw[->, thick] (state) -- (plan);
      \draw[->, dashed, thick] (plan.south) to[out=-70,in=-20] (state.south);
    \end{tikzpicture}
    \caption{Recreated teaching figure: two predictive emphases.}
  \end{figure}
\end{frame}

\begin{frame}{How to Read the Figure}
  The top path asks:
  \begin{itemize}
    \item what comes next,
    \item how should the sequence continue,
    \item how do we maximize fluent continuation?
  \end{itemize}

  The bottom path asks:
  \begin{itemize}
    \item what state of the world am I in,
    \item what happens if I act,
    \item which imagined future should I prefer?
  \end{itemize}
\end{frame}

\begin{frame}{What World Models Buy You}
  World models are one of the strongest candidates for progress beyond language-only systems because they target:
  \begin{itemize}
    \item prediction,
    \item abstraction,
    \item planning,
    \item action.
  \end{itemize}

  They are especially compelling where:
  \begin{itemize}
    \item control matters,
    \item consequences matter,
    \item interventions change future observations.
  \end{itemize}
\end{frame}

\begin{frame}{Why the Evidence Is Still Narrower Than the Hype}
  The strongest evidence for world models is concentrated in:
  \begin{itemize}
    \item reinforcement learning,
    \item robotics,
    \item video representation,
    \item planning-heavy tasks.
  \end{itemize}

  \begin{warningbox}{Do not overclaim}
    It is not yet well supported to say that JEPA-style methods have already superseded autoregressive modeling across all of AI.
  \end{warningbox}

  \source{\href{https://arxiv.org/abs/2601.00844}{Destrade et al. (2026): Value-guided JEPA Planning}}
\end{frame}

\begin{frame}{V-JEPA 2 and the Move Toward Robotics}
  V-JEPA 2 is especially notable because it pushes from large-scale video prediction toward robotic planning with a small amount of robot-interaction data.

  That matters because it shows a plausible route from:
  \begin{itemize}
    \item predictive latent representation,
    \item to planning,
    \item to embodied deployment.
  \end{itemize}

  \source{\href{https://arxiv.org/abs/2506.09985}{Assran et al. (2025): V-JEPA 2}}
\end{frame}

\section{Embodied Intelligence}

\begin{frame}{What Embodiment Actually Means}
  Embodiment does not simply mean ``having a camera'' or ``being multimodal.''

  \begin{conceptbox}{Definition}
    Embodiment means intelligence coupled to a body that acts in an environment, changes future observations through action, and must manage the sensorimotor consequences of its own behavior.
  \end{conceptbox}

  \source{\href{https://people.csail.mit.edu/brooks/idocs/elephants.pdf}{Brooks (1990): Elephants Don't Play Chess}}
\end{frame}

\begin{frame}{Why Physical Deployment Changes the Problem}
  Some useful AI systems do not need bodies in the strong sense:
  \begin{itemize}
    \item theorem assistants,
    \item schedulers,
    \item CAD copilots,
    \item document-review systems.
  \end{itemize}

  But once the world can push back through friction, delay, occlusion, contact, or damage, text-only competence becomes a weak proxy for intelligence.
\end{frame}

\begin{frame}{Why Robotics Put Embodiment Back at the Center}
  Modern robotics research explains why embodiment returned to the center of AI discussion:
  \begin{itemize}
    \item RT-1 and RT-2 showed broad robot control with transformer-style architectures,
    \item Open X-Embodiment demonstrated positive transfer across robots and institutions,
    \item OpenVLA and Octo pushed the generalist robot policy idea further.
  \end{itemize}

  \source{\href{https://arxiv.org/abs/2307.15818}{Brohan et al. (RT-2)}; \href{https://arxiv.org/abs/2310.08864}{Open X-Embodiment Collaboration}; \href{https://arxiv.org/abs/2406.09246}{OpenVLA}}
\end{frame}

\begin{frame}{From Language Interface to Physical Agency}
  \begin{figure}
    \centering
    \begin{tikzpicture}[>=Stealth, font=\small, node distance=0.45cm, scale=0.82, transform shape]
      \tikzset{stage/.style={draw, rounded corners, minimum width=2.55cm, minimum height=0.95cm, align=center}}
      \node[stage, fill=blue!8] (text) {Language-only systems\\chat, coding, retrieval};
      \node[stage, fill=green!10, right=of text] (tools) {Tool-using systems\\search, code execution, workflows};
      \node[stage, fill=orange!12, right=of tools] (models) {World-aware agents\\prediction, planning, simulation};
      \node[stage, fill=red!9, right=of models] (emb) {Embodied systems\\robotics, navigation, manipulation};
      \draw[->, thick] (text) -- (tools);
      \draw[->, thick] (tools) -- (models);
      \draw[->, thick] (models) -- (emb);
      \node[align=center, below=0.55cm of tools] {Increasing need for persistent state,\\causal prediction, and physical feedback};
    \end{tikzpicture}
    \caption{Recreated teaching figure: from interface intelligence to physical agency.}
  \end{figure}
\end{frame}

\begin{frame}{Biological Analogies: Use Carefully}
  The biological examples are useful pedagogically if interpreted carefully:
  \begin{itemize}
    \item paramecium and amoebae show adaptive environment-coupled behavior below brains,
    \item crows and dolphins show that sophisticated cognition is not confined to primates.
  \end{itemize}

  The lesson is not that robots are about to become animals. The lesson is that intelligence is biologically plural and strongly shaped by bodies and niches.
\end{frame}

\begin{frame}{What Embodiment Does \emph{Not} Imply}
  \begin{warningbox}{Do not universalize from robotics}
    Embodiment is necessary for robust physical agency, but it is not obviously necessary for every economically useful form of AI competence.
  \end{warningbox}

  The stronger claim is domain-specific:
  \begin{itemize}
    \item when the world can push back through physics,
    \item embodiment becomes increasingly important,
    \item but not every useful AI system must become embodied.
  \end{itemize}
\end{frame}

\begin{frame}{Engineering Takeaway on Embodiment}
  If the task is:
  \begin{itemize}
    \item closed-loop,
    \item safety-critical,
    \item contact-rich,
    \item irreversible,
  \end{itemize}
  then embodiment becomes a first-order engineering question rather than a philosophical embellishment.

  \begin{examplebox}{Examples}
    maintenance robotics, navigation, manipulation, human-robot interaction, industrial autonomy.
  \end{examplebox}
\end{frame}

\section{AGI, ASI, and Oversight}

\begin{frame}{Why AGI and ASI Are Slippery Terms}
  Artificial general intelligence and artificial superintelligence are among the least disciplined terms in public AI discourse.

  Problems:
  \begin{itemize}
    \item no consensus operational benchmark,
    \item performance can be confused with intelligence,
    \item impressive task demonstrations can be over-generalized into species-level narratives.
  \end{itemize}

  \source{\href{https://arxiv.org/abs/2311.02462}{Morris et al. (Levels of AGI)}; \href{https://arxiv.org/abs/1911.01547}{Chollet (2019)}}
\end{frame}

\begin{frame}{Why Capability Talk Jumps Too Fast}
  The literature discussed here does \emph{not} demonstrate a system with all of the following together:
  \begin{itemize}
    \item broad autonomous competence across open-ended environments,
    \item stable self-improvement under weak oversight,
    \item robust goal alignment under superhuman capability,
    \item reliable control under distribution shift.
  \end{itemize}

  Those are separate requirements, and the evidence does not yet show them together.
\end{frame}

\begin{frame}{What Serious Near-Term Work Actually Studies}
  The most serious work near this area is not proof of ASI. It is \textbf{oversight research} motivated by the possibility that future systems may outstrip direct human inspection.

  Examples:
  \begin{itemize}
    \item Constitutional AI,
    \item weak-to-strong generalization,
    \item debate and consultancy protocols.
  \end{itemize}

  \source{\href{https://arxiv.org/abs/2212.08073}{Bai et al. (Constitutional AI)}; \href{https://arxiv.org/abs/2312.09390}{Burns et al. (Weak-to-Strong)}}
\end{frame}

\begin{frame}{Engineering Standard of Evidence}
  \begin{conceptbox}{A practical rule}
    A strong future-AI claim should be evaluated the same way any engineering claim is evaluated: what was demonstrated, under what conditions, with what controls, with what failure modes, and with what verification mechanism?
  \end{conceptbox}

  This standard is deliberately stricter than hype-cycle rhetoric.
\end{frame}

\begin{frame}{What the Literature Does Not Yet Justify}
  The responsible teaching stance is neither ridicule nor prophecy.

  It is to say clearly:
  \begin{itemize}
    \item the evidentiary burden is extreme,
    \item current results are interesting but partial,
    \item oversight and governance are more concrete than superintelligence storytelling.
  \end{itemize}
\end{frame}

\section{What Engineers Should Conclude}

\begin{frame}{The Likely Next Decade: Diversification of the Stack}
  From an engineering point of view, the likely next decade is not a clean handoff from LLMs to one replacement architecture.

  It is a diversification of the stack:
  \begin{itemize}
    \item foundation models as interface layers,
    \item neuro-symbolic methods where explicit structure matters,
    \item world models where prediction and planning matter,
    \item embodiment where physical deployment matters.
  \end{itemize}
\end{frame}

\begin{frame}{Distinction Sheet I}
  \small
  \begin{tabular}{p{0.22\textwidth}p{0.34\textwidth}p{0.34\textwidth}}
    \toprule
    \textbf{Theme} & \textbf{Established} & \textbf{Strong hypothesis} \\
    \midrule
    Current LLM systems & Powerful for language, code, retrieval, and tool-using workflows & Better scaffolding can keep extending usefulness in many professional settings \\
    Neuro-symbolic AI & Hybrid systems can improve structured reasoning and explicit constraint handling & Constraint-aware hybrids will become more important in engineering-grade AI \\
    World models & Learned latent dynamics help control and planning in RL and robotics & Predictive state models are necessary for robust physical agency \\
    \bottomrule
  \end{tabular}
\end{frame}

\begin{frame}{Distinction Sheet II}
  \small
  \begin{tabular}{p{0.22\textwidth}p{0.34\textwidth}p{0.34\textwidth}}
    \toprule
    \textbf{Theme} & \textbf{Established} & \textbf{Speculative} \\
    \midrule
    Embodiment & Physical interaction changes the competence required & Every useful future AI system must be embodied \\
    AGI / ASI & There is no consensus operational benchmark & Near-term ASI is already an established engineering forecast \\
    \bottomrule
  \end{tabular}
\end{frame}

\begin{frame}{The Less Credible Forecast}
  The less credible forecast is the one preferred by hype cycles:
  \begin{itemize}
    \item one architecture becomes universally intelligent,
    \item human oversight fades away,
    \item engineering collapses into prompting.
  \end{itemize}

  \begin{warningbox}{Why this is weak}
    The literature supports a more demanding picture in which future AI systems become more heterogeneous, more infrastructure-like, and more dependent on careful interface design, deployment boundaries, and verification.
  \end{warningbox}
\end{frame}

\begin{frame}{Final Message}
  The frontier after current LLM systems is real.

  \begin{itemize}
    \item Neuro-symbolic AI says that intelligence without explicit structure is often not enough.
    \item World-model research says that prediction and planning over latent state may matter more than surface continuation when action matters.
    \item Embodiment says that intelligence tested only in language is not the same as intelligence tested in the world.
    \item AGI/ASI discourse says that capability talk without measurement and oversight is mostly theater.
  \end{itemize}
\end{frame}

\begin{frame}{Engineers Stay Above AI}
  \begin{conceptbox}{Chapter conclusion}
    Regardless of paradigm, engineers stay above AI because engineers still define the task, the constraints, the acceptable evidence, the failure boundaries, and the accountability chain.
  \end{conceptbox}

  That is the design principle that should survive every architecture transition discussed in this course.
\end{frame}

\end{document}
